\chapter{Speed costs tilt, and tilt costs view}
\label{chap:steady}

\begin{lemma}[Body-frame facts]
\label{lem:frame}
\lean{VisionTracking.QuadTraj.b₃_unit, VisionTracking.QuadTraj.b₁₃_orth,
  VisionTracking.QuadTraj.thrust_nonneg, VisionTracking.norm_cross_of_orthonormal,
  VisionTracking.inner_cross_left, VisionTracking.inner_cross_right}
\leanok
\uses{def:quad-traj}
In every state, $b_3=b_1\times b_2$ is a unit vector orthogonal to $b_1$.
\end{lemma}

\begin{theorem}[Pointing bound for every quadrotor state]
\label{thm:quad-pointing}
\lean{VisionTracking.QuadTraj.abs_sub_elevation_le_angle,
  VisionTracking.QuadTraj.elevation_mem_of_sees}
\leanok
\uses{lem:frame, lem:orthonormal-bound}
In every state, the camera misses a unit direction $u$ by at least
$|\eta-e(b_3,u)|$. So if it sees $u$, then
$e(b_3,u)\in[\eta_{\mathrm{lo}}-\alpha,\eta_{\mathrm{hi}}+\alpha]$.
\end{theorem}

\begin{definition}[Chase tilt]
\label{def:chase-tilt}
\lean{VisionTracking.chaseTilt, VisionTracking.steadyAxis}
\leanok
$\theta(v)=\arctan(c_Dv^2/g)$, and the steady axis
$\cos\theta\,\up+\sin\theta\,\hat w$.
\end{definition}

\begin{theorem}[Steady flight forces the tilt]
\label{thm:forced-tilt}
\lean{VisionTracking.b₃_eq_steadyAxis, VisionTracking.eq_of_smul_eq_smul_of_unit}
\leanok
\uses{def:chase-tilt, lem:frame}
If the velocity is horizontal, nonzero and unchanged over a step, then whatever the
motors do, $b_3$ is the steady axis leaning by $\theta(\norm{v})$ towards the
direction of travel.
\end{theorem}

\begin{theorem}[The steady-chase window]
\label{thm:window}
\lean{VisionTracking.steady_chase_window, VisionTracking.steady_visible_iff,
  VisionTracking.inplane_visible_iff, VisionTracking.inner_steadyAxis,
  VisionTracking.norm_steadyAxis}
\leanok
\uses{thm:forced-tilt, thm:quad-pointing}
In steady flight, a target seen ahead in the plane of travel at elevation $\lambda$
must satisfy $\eta_{\mathrm{lo}}-\alpha\le\theta(v)+\lambda\le\eta_{\mathrm{hi}}+\alpha$.
\end{theorem}

\begin{corollary}[Trackable speeds depend on the altitude offset]
\label{cor:window-speed}
\lean{VisionTracking.chaseTilt_add_le_iff, VisionTracking.le_chaseTilt_add_iff,
  VisionTracking.arctan_le_iff_le_tan, VisionTracking.chaseTilt_nonneg,
  VisionTracking.chaseTilt_lt}
\leanok
\uses{thm:window}
$c_Dv^2\le g\tan(\eta_{\mathrm{hi}}+\alpha-\lambda)$ and
$g\tan(\eta_{\mathrm{lo}}-\alpha-\lambda)\le c_Dv^2$. From directly behind at the same
altitude, the top speed is $\sqrt{g\tan(\eta_{\mathrm{hi}}+\alpha)/c_D}$. Flying
above the target raises it; flying below lowers it.
\end{corollary}

\section{Lifting plans to motors and attitude}

\begin{theorem}[Motor allocation]
\label{thm:motors}
\lean{VisionTracking.exists_motors_of_wrench, VisionTracking.QuadParams.mix,
  VisionTracking.totalThrust_mix, VisionTracking.torque_mix}
\leanok
\uses{def:mixing}
The mixer has an explicit inverse. A wrench $(F,\tau)$ whose per-axis demands fit
within the thrust margins is produced by motors in $[0,f_{\mathrm{motor}}]$.
\end{theorem}

\begin{theorem}[Rodrigues is a proper rotation]
\label{thm:rotation}
\lean{VisionTracking.inner_rotate, VisionTracking.norm_rotate, VisionTracking.rotate_cross,
  VisionTracking.angle_rotate_le, VisionTracking.rodr, VisionTracking.inner_rodr,
  VisionTracking.rodr_cross}
\leanok
\uses{def:rotate}
$\mathrm{rotate}(\omega)$ preserves inner and cross products, so attitude updates keep
the frame right-handed orthonormal. It turns any vector by at most $\norm{\omega}$.
The polynomial identities are proved by \texttt{linear\_combination} with
Gröbner-basis certificates.
\end{theorem}

\begin{theorem}[Axis-angle for nearby frames]
\label{lem:axis-angle}
\lean{VisionTracking.exists_rotate_frame, VisionTracking.frameTrace, VisionTracking.rotAngle,
  VisionTracking.frameTrace_gt_of_angle_le, VisionTracking.rotAngle_le,
  VisionTracking.core_rotate, VisionTracking.ofF, VisionTracking.toF,
  VisionTracking.cross_ofF, VisionTracking.ofF_toF, VisionTracking.rotate_ofF}
\leanok
\uses{thm:rotation}
Two right-handed orthonormal frames whose trace $t=\sum_i\inner{a_i}{c_i}$ exceeds $-1$
differ by a single $\mathrm{rotate}(\omega)$, with $\norm{\omega}=\arccos((t-1)/2)$. If
every frame vector moves by at most $\varepsilon\le\pi/2$, then the angle is at most
$2\varepsilon$.
\end{theorem}

\begin{proof}
\leanok
Map to the coordinates of the first frame by an isometry that preserves cross
products. There, $\sigma=\tfrac12\sum_i e_i\times c_i$ satisfies
$(1+\cos\theta)(c_i-\cos\theta\,e_i-\sigma\times e_i)=\inner{\sigma}{e_i}\sigma$ and
$\norm{\sigma}^2=1-\cos^2\theta$. These are polynomial identities modulo
orthonormality, certified by \texttt{linear\_combination}. For the angle bound,
$t\ge3\cos\varepsilon$ gives $(t-1)/2\ge\cos2\varepsilon$.
\end{proof}

\begin{theorem}[Lifting kinematic plans]
\label{thm:lifting}
\lean{VisionTracking.QuadTraj.ofPlan, VisionTracking.KinPlan, VisionTracking.MixerOK,
  VisionTracking.KinPlan.thrust, VisionTracking.KinPlan.rate, VisionTracking.KinPlan.torque,
  VisionTracking.KinPlan.rotVec, VisionTracking.KinPlan.rotVec_spec,
  VisionTracking.KinPlan.smul_worldRate_rate, VisionTracking.inertiaInv_inertia,
  VisionTracking.ofPlan_x, VisionTracking.ofPlan_b₁, VisionTracking.ofPlan_b₃,
  VisionTracking.ofPlan_η}
\leanok
\uses{lem:axis-angle, thm:motors}
Take a plan of positions, velocities, body frames (consecutive ones within one
rotation) and hinge angles. Suppose the required thrust
$T_k=\Delta v_k/\delta+g\up+c_D\norm{v_k}v_k$ points along $b_3$, and the wrench
$(m\norm{T_k},\tau_k)$ with $\tau_k=J\Delta\Omega_k/\delta+\Omega_k\times J\Omega_k$ fits
the mixer. Suppose also that tilt, body rates, speed and hinge are within limits.
Then the plan is a valid rigid-body trajectory.
\end{theorem}

\begin{lemma}[Regularity of normalisation and of the downhill axis]
\label{lem:regularity}
\lean{VisionTracking.nrm, VisionTracking.norm_nrm_sub_le, VisionTracking.norm_nrm_second_le,
  VisionTracking.norm_add_le_quad, VisionTracking.second_scalar, VisionTracking.inv_sub_mul_le,
  VisionTracking.downhill, VisionTracking.downhillRaw, VisionTracking.downhill_unit,
  VisionTracking.downhill_orth, VisionTracking.downhill_aligned,
  VisionTracking.norm_downhillRaw_sub_le, VisionTracking.norm_downhillRaw_second_le}
\leanok
On vectors of length at least $m$:
$\norm{\hat a-\hat b}\le2\norm{a-b}/m$ and
$\norm{\hat c-2\hat b+\hat a}\le2\norm{c-2b+a}/m+3(\norm{c-b}^2+\norm{b-a}^2)/m^2$.
The downhill axis $\widehat{-\up+\inner{b}{\up}b}$ is a unit vector orthogonal to $b$,
aligned with the rotor-plane part of $-\up$. Its unnormalised form has first
differences at most $2\norm{\Delta b}$ and second differences at most
$2\norm{\Delta^2b}+\norm{\Delta b}(\norm{\Delta b}+\norm{\Delta b'})$.
\end{lemma}

\begin{theorem}[Body rates from frame differences]
\label{thm:frame-rate}
\lean{VisionTracking.KinPlan.norm_rate_le, VisionTracking.KinPlan.norm_rate_sub_le,
  VisionTracking.frameTrace_ge_of_step, VisionTracking.psiVec, VisionTracking.core_rotate_psi,
  VisionTracking.frameRotVec, VisionTracking.frameRotVec_spec, VisionTracking.norm_psiVec_le,
  VisionTracking.norm_psiVec_sub_le, VisionTracking.norm_sigmaVec_sub_le,
  VisionTracking.arcsin_sub_le, VisionTracking.KinPlan.rate_eq}
\leanok
\uses{lem:axis-angle, thm:lifting}
For rotations of at most $\pi/2$, the axis-angle vector is
$(\arcsin\norm\sigma/\norm\sigma)\sigma$, which is $3/\sqrt{1-\rho^2}$-Lipschitz on
$\norm\sigma\le\rho$. If the frame vectors move by at most $D_1$ per step with second
differences at most $D_2$, and $2D_1<1$, then
$\norm{\Omega_k}\le2D_1/(\delta\sqrt{1-4D_1^2})$ and
$\norm{\Omega_{k+1}-\Omega_k}\le18(D_2+D_1^2)/(\delta\sqrt{1-4D_1^2})$.
\end{theorem}

\section{Steady formation: exact characterization}

\begin{theorem}[Steady formation iff window]
\label{thm:steady-formation}
\lean{VisionTracking.steady_formation_iff, VisionTracking.levelTarget,
  VisionTracking.steady_thrust, VisionTracking.exists_aligned,
  VisionTracking.cross_cross_self}
\leanok
\uses{thm:window, thm:motors, thm:rotation, thm:pointing}
Against a target flying straight and level at speed $v$, a valid rigid-body
trajectory holding the target at elevation $\lambda$ and range $r\le\ddet$ in
steady formation exists \emph{iff}
$\eta_{\mathrm{lo}}-\alpha\le\theta(v)+\lambda\le\eta_{\mathrm{hi}}+\alpha$. This
assumes the four motors can produce $m\sqrt{g^2+(c_Dv^2)^2}$ and that
$\theta(v)\le\theta_{\max}$.
\end{theorem}

\begin{proof}
\leanok
$\Rightarrow$ is Theorem~\ref{thm:window}. $\Leftarrow$: constant velocity; $b_3$
on the steady axis; $b_1$ in the rotor plane towards the line of sight;
$b_2=b_3\times b_1$; zero body rates; four equal motor thrusts; the hinge clamped to
the elevation.
\end{proof}
